\documentclass{article}
\usepackage{geometry}
\geometry{margin=1.5cm, vmargin={0pt,1cm}}
\setlength{\topmargin}{-1cm}
\setlength{\headheight}{12.5pt}
\setlength{\paperheight}{29.7cm}
\setlength{\textheight}{25.3cm}

% useful packages.
\usepackage[UTF8]{ctex}
\usepackage{fancyhdr}
\usepackage{layout}
\usepackage{luacode}
\usepackage{listings}
\usepackage{xcolor}

% Configure listings for Python code highlighting
\lstset{
  language=Python,
  basicstyle=\ttfamily\small,
  keywordstyle=\bfseries\color{blue},
  commentstyle=\color{gray},
  stringstyle=\color{red},
  breaklines=true,
  showstringspaces=false,
  tabsize=4,
  frame=single,
  framesep=5pt,
  rulecolor=\color{gray},
  backgroundcolor=\color{white},
  captionpos=b,
  numberstyle=\tiny\color{gray},
  numbers=left,
  stepnumber=5,
}

\lstnewenvironment{plaintext}[1][]{
  \lstset{
    language=none,
    basicstyle=\ttfamily,
    keywordstyle=,
    commentstyle=,
    stringstyle=,
    numbers=none,
    frame=none,
    backgroundcolor=\color{white},
    #1                % 允许用户自行覆盖默认设置
  }
}{}

% import common macros
% % % % % % % % % % % % % % % % % % % % % % % % % % % % % % % %
% Common Macros for LaTeX
% % % % % % % % % % % % % % % % % % % % % % % % % % % % % % % %

% Useful packages
\usepackage{amsfonts}
\usepackage{amsmath}
\usepackage{amssymb}
\usepackage{amsthm}
\usepackage{enumerate}
\usepackage{graphicx}
\usepackage{multicol}
\usepackage[ruled,linesnumbered]{algorithm2e}

% Math operators and common symbols
\newcommand{\dif}{\mathrm{d}}
\newcommand{\innerProd}[1]{\left\langle #1 \right\rangle}
\newcommand{\difFrac}[2]{\frac{\dif #1}{\dif #2}}
\newcommand{\pdfFrac}[2]{\frac{\partial #1}{\partial #2}}
\newcommand{\T}{\mathrm{T}}
\newcommand{\norm}[1]{\left\| #1 \right\|}
\newcommand{\abs}[1]{\left| #1 \right|}

% Vector and Matrix shortcuts
\newcommand{\vecTwo}[2]{
  \begin{bmatrix}
    #1 \\
    #2
\end{bmatrix}}
\newcommand{\vecThree}[3]{
  \begin{bmatrix}
    #1 \\
    #2 \\
    #3
\end{bmatrix}}
\newcommand{\vecTwoH}[2]{
  \begin{bmatrix}
    #1 & #2
  \end{bmatrix}
}
\newcommand{\vecThreeH}[3]{
  \begin{bmatrix}
    #1 & #2 & #3
  \end{bmatrix}
}

% Common fields
\newcommand{\R}{\mathbb{R}}
\newcommand{\C}{\mathbb{C}}
\newcommand{\Z}{\mathbb{Z}}
\newcommand{\N}{\mathbb{N}}

% Solutions and proofs
\newenvironment{solution}{
\begin{proof}[解]}{
\end{proof}}

% Utility to get environment variables (requires lualatex)
\newcommand{\getEnv}[2]{\directlua{tex.print(os.getenv("#1") or "#2")}}


% Name and uID
\newcommand{\name}{\getEnv{NAME}{NAME}}
\newcommand{\uid}{\getEnv{UID}{UID}}
\newcommand{\course}{数值代数}
\newcommand{\hwNum}{4}

\begin{document}

\pagestyle{fancy}
\fancyhead{}
\lhead{\name (\uid)}
\chead{\course \#\hwNum}
\rhead{\today}

\section{题目 1：向量范数}

设 $\|\cdot\|_p: \mathbb{R}^n \to \mathbb{R}$ 是由
\[
  \|x\|_p = \left( |x_1|^p + \cdots + |x_n|^p \right)^{1/p}, \quad x
  \in \mathbb{R}^n, \ p \geq 1
\]
定义的。证明：
\begin{enumerate}
  \item $\|\cdot\|_p$ 是一个向量范数。
  \item $\|x\|_\infty = \max_{1 \leq i \leq n} |x_i|$。
\end{enumerate}

\begin{solution}
  \textbf{Step 1. 证明 H\"older 不等式}

  首先证明 Young 不等式：对于任意 $a, b \geq 0$ 和 $p, q > 1$ 满足 $\frac{1}{p} +
  \frac{1}{q} = 1$，有
  \[
    ab \leq \frac{a^p}{p} + \frac{b^q}{q}.
  \]

  考虑对数 $\ln(ab) = \ln a + \ln b = \frac{1}{p} \ln a^p + \frac{1}{q} \ln b^q$。
  于是由 Jensen 不等式可得 $\ln(ab) \leq \ln\left( \frac{a^p}{p} +
  \frac{b^q}{q} \right)$，从而得到 Young 不等式。

  然后考虑 H\"older 不等式：对于任意 $x, y \in \mathbb{R}^n$ 和 $p, q > 1$ 满足 $\frac{1}{p} +
  \frac{1}{q} = 1$，有

  \[
    \sum_{i=1}^{n} |x_i y_i| \leq \left( \sum_{i=1}^{n} |x_i|^p \right)^{1/p}
    \left( \sum_{i=1}^{n} |y_i|^q \right)^{1/q}.
  \]

  首先考虑归一化。令 $A = \left( \sum_{i=1}^{n} |x_i|^p \right)^{1/p}$ 和 $B =
  \left( \sum_{i=1}^{n} |y_i|^q \right)^{1/q}$，则 $A, B > 0$。于是令
  $\hat{a}_i = \frac{|x_i|}{A}$ 和 $\hat{b}_i = \frac{|y_i|}{B}$，则
  $\sum_{i=1}^{n} \hat{a}_i^p = \sum_{i=1}^{n} \hat{b}_i^q = 1$。由 Young 不等式可得
  \[
    \sum_{i=1}^{n} \hat{a}_i \hat{b}_i \leq \sum_{i=1}^{n} \left(
      \frac{\hat{a}_i^p}{p} +
    \frac{\hat{b}_i^q}{q} \right) = 1.
  \]
  从而
  \[
    \sum_{i=1}^{n} |x_i y_i| = AB \sum_{i=1}^{n} \hat{a}_i \hat{b}_i \leq AB =
    \left( \sum_{i=1}^{n} |x_i|^p \right)^{1/p} \left( \sum_{i=1}^{n} |y_i|^q
    \right)^{1/q}.
  \]
  从而得到 H\"older 不等式。

  \textbf{Step 2. 证明 $\|\cdot\|_p$ 是一个向量范数}

  首先正定性与线性性显然，只需证明三角不等式。

  考虑 $|x_i + y_i|^p = |x_i + y_i| |x_i + y_i|^{p-1} \leq (|x_i| +
  |y_i|) |x_i + y_i|^{p-1}$ ，求和，得到：
  \[
    \sum_{i=1}^{n} |x_i + y_i|^p \leq \sum_{i=1}^{n} |x_i| |x_i + y_i|^{p-1} +
    \sum_{i=1}^{n} |y_i| |x_i + y_i|^{p-1}.
  \]

  考虑 $\sum |x_i||x_i+y_i|^{p-1}$ 的项，令 $a_i = |x_i|$ 和 $b_i = |x_i +
  y_i|^{p-1}$，并令 $q = \frac{p}{p-1}$，则 $a_i^p = |x_i|^p$ 和 $b_i^q =
  |x_i + y_i|^p$。由 H\"older 不等式可得

  \[
    \sum_{i=1}^{n} |x_i||x_i+y_i|^{p-1} \leq \|x\|_p \cdot \|x+y\|_q^{p/q}
  \]

  对 $\sum |y_i||x_i+y_i|^{p-1}$ 的项同样处理，得到

  \[
    \sum_{i=1}^{n} |y_i||x_i+y_i|^{p-1} \leq \|y\|_p \cdot \|x+y\|_q^{p/q}.
  \]

  于是有：

  \[
    \sum_{i=1}^{n} |x_i + y_i|^p \leq (\|x\|_p + \|y\|_p) \cdot \|x+y\|_q^{p/q}.
  \]

  这里又有：$p/q = p - 1$，所以 $\|x+y\|_q^{p/q} = \|x+y\|_p^{p-1}$。因此：

  \[
    \|x+y\|_p^p \leq (\|x\|_p + \|y\|_p) \cdot \|x+y\|_p^{p-1},
    \Rightarrow \|x+y\|_p \leq \|x\|_p + \|y\|_p.
  \]

  从而得证。

  \textbf{Step 3. 证明 $\|x\|_\infty = \max_{1 \leq i \leq n} |x_i|$}

  不妨设 $\max_{1 \leq i \leq n} |x_i| = |x_1|$，则 $|x_i| \leq |x_1|$ 对任意
  $i$ 都成立，因此 $\|x\|_p \leq (n |x_1|^p)^{1/p} = n^{1/p} |x_1|$。当 $p
  \to \infty$ 时，$n^{1/p} \to 1$，因此 $\|x\|_\infty \leq \max_{1 \leq i
  \leq n} |x_i|$。

  另一方面，$|x_1| \leq \|x\|_p$ 对任意 $p$ 都成立，因此 $\max_{1 \leq i \leq n}
  |x_i| \leq \|x\|_\infty$。

  从而得证。

\end{solution}

\section{题目 2：向量范数的等价性}

证明 $\|\cdot\|_1$, $\|\cdot\|_2$, $\|\cdot\|_\infty$ 有如下的等价性：
\begin{enumerate}
  \item $\|x\|_2 \leq \|x\|_1 \leq \sqrt{n}\|x\|_2$。
  \item $\|x\|_\infty \leq \|x\|_2 \leq \sqrt{n}\|x\|_\infty$。
  \item $\|x\|_\infty \leq \|x\|_1 \leq n\|x\|_\infty$。
\end{enumerate}

\begin{solution}
  \subsection*{(1) $L_1$ 范数与 $L_2$ 范数的等价性}

  首先证明 $\|x\|_2 \leq \|x\|_1$。由于
  \[
    \|x\|_1^2 = \left( \sum_{i=1}^{n} |x_i| \right)^2 =
    \sum_{i=1}^{n} |x_i|^2 + 2 \sum_{1 \leq i < j \leq n} |x_i||x_j|
    \geq \sum_{i=1}^{n} |x_i|^2 = \|x\|_2^2,
  \]
  因此 $\|x\|_2 \leq \|x\|_1$。

  然后证明 $\|x\|_1 \leq \sqrt{n} \|x\|_2$。由 Cauchy-Schwarz 不等式：
  \[
    \|x\|_1 = \sum_{i=1}^{n} |x_i| \cdot 1 \leq \left( \sum_{i=1}^{n}
    |x_i|^2 \right)^{1/2} \left( \sum_{i=1}^{n} 1^2 \right)^{1/2} =
    \sqrt{n} \|x\|_2.
  \]

  综上，$\|x\|_2 \leq \|x\|_1 \leq \sqrt{n} \|x\|_2$。

  \subsection*{(2) $L_2$ 范数与 $L_\infty$ 范数的等价性}

  设 $|x_k| = \|x\|_\infty = \max_{1 \leq i \leq n} |x_i|$，则
  \[
    \|x\|_2 = \left( \sum_{i=1}^{n} |x_i|^2 \right)^{1/2} \geq \left(
    |x_k|^2 \right)^{1/2} = |x_k| = \|x\|_\infty.
  \]
  这给出了 $\|x\|_\infty \leq \|x\|_2$。

  另一方面，
  \[
    \|x\|_2 = \left( \sum_{i=1}^{n} |x_i|^2 \right)^{1/2} \leq \left(
    \sum_{i=1}^{n} \|x\|_\infty^2 \right)^{1/2} = \left( n
    \|x\|_\infty^2 \right)^{1/2} = \sqrt{n} \|x\|_\infty.
  \]

  综上，$\|x\|_\infty \leq \|x\|_2 \leq \sqrt{n} \|x\|_\infty$。

  \subsection*{(3) $L_1$ 范数与 $L_\infty$ 范数的等价性}

  设 $|x_k| = \|x\|_\infty$，则
  \[
    \|x\|_1 = \sum_{i=1}^{n} |x_i| \geq |x_k| = \|x\|_\infty.
  \]

  另一方面，
  \[
    \|x\|_1 = \sum_{i=1}^{n} |x_i| \leq \sum_{i=1}^{n} \|x\|_\infty =
    n \|x\|_\infty.
  \]

  综上，$\|x\|_\infty \leq \|x\|_1 \leq n \|x\|_\infty$。

\end{solution}

\section{题目 3：Frobenius 范数}

设 $\|\cdot\|_F: \mathbb{R}^{n \times n} \to \mathbb{R}$ 是由
\[
  \|A\|_F = \left( \sum_{i,j=1}^{n} |a_{ij}|^2 \right)^{1/2}, \quad A
  = (a_{ij}) \in \mathbb{R}^{n \times n}
\]
定义的。证明 $\|\cdot\|_F$ 是一个矩阵范数。（注：该范数也是一个常用且易于计算的矩阵范数，称为 Frobenius 范数。）

\begin{solution}
  正定性、线性性与三角不等式由 $L_2$ 范数的性质易得，因此只需证明相容性，即 $\|AB\|_F \leq \|A\|_F \|B\|_F$。

  令 $C = AB$，则 $c_{ij} = \sum_{k=1}^{n} a_{ik} b_{kj}$。由 Cauchy-Schwarz 不等式：
  \[
    |c_{ij}|^2 = \left| \sum_{k=1}^{n} a_{ik} b_{kj} \right|^2 \leq
    \left( \sum_{k=1}^{n} |a_{ik}|^2 \right) \left( \sum_{k=1}^{n}
    |b_{kj}|^2 \right).
  \]

  对所有 $i, j$ 求和：
  \[
    \|AB\|_F^2 = \sum_{i,j=1}^{n} |c_{ij}|^2 \leq \sum_{i,j=1}^{n}
    \left( \sum_{k=1}^{n} a_{ik}^2 \right) \left( \sum_{k=1}^{n}
    b_{kj}^2 \right).
  \]

  注意到 $\sum_{k=1}^{n} a_{ik}^2$ 只与 $i$ 有关，$\sum_{k=1}^{n} b_{kj}^2$
  只与 $j$ 有关，因此：
  \[
    \|AB\|_F^2 \leq \left( \sum_{i=1}^{n} \sum_{k=1}^{n} a_{ik}^2
    \right) \left( \sum_{j=1}^{n} \sum_{k=1}^{n} b_{kj}^2 \right) =
    \|A\|_F^2 \|B\|_F^2.
  \]

  从而 $\|AB\|_F \leq \|A\|_F \|B\|_F$，相容性得证。

\end{solution}

\section{题目 4：矩阵范数的不等式}

设 $A, B \in \mathbb{R}^{n \times n}$，证明：
\begin{enumerate}
  \item $\max_{1 \leq i,j \leq n} |a_{ij}| \leq \|A\|_2 \leq n
    \max_{1 \leq i,j \leq n} |a_{ij}|$。
  \item $\frac{1}{\sqrt{n}} \|A\|_\infty \leq \|A\|_2 \leq
    \sqrt{n}\|A\|_\infty$。
  \item $\frac{1}{\sqrt{n}} \|A\|_1 \leq \|A\|_2 \leq \sqrt{n}\|A\|_1$。
  \item $\|A\|_2 \leq \|A\|_F \leq \sqrt{n}\|A\|_2$。
  \item $\|AB\|_F \leq \|A\|_2\|B\|_F$。
  \item $\|AB\|_F \leq \|A\|_F\|B\|_2$。
\end{enumerate}

\begin{solution}
  \subsection*{(1) $\max_{i,j} |a_{ij}| \leq \|A\|_2 \leq n
  \max_{i,j} |a_{ij}|$}

  设 $\max_{i,j} |a_{ij}| = |a_{kl}|$。取 $e_k$ 为第 $k$ 个单位向量，则
  \[
    \|A\|_2 \geq \|Ae_k\|_2 = \sqrt{\sum_{i=1}^n a_{ik}^2} \geq
    |a_{kl}| = \max_{i,j} |a_{ij}|.
  \]

  另一方面，设 $\|x\|_2 = 1$，则
  \[
    \|Ax\|_2^2 = \sum_{i=1}^n \left( \sum_{j=1}^n a_{ij} x_j
    \right)^2 \leq \sum_{i=1}^n \left( n \max_{i,j} |a_{ij}|^2 \cdot
    \sum_{j=1}^n x_j^2 \right) = n^2 \max_{i,j} |a_{ij}|^2.
  \]
  因此 $\|A\|_2 \leq n \max_{i,j} |a_{ij}|$。

  \subsection*{(2) $\frac{1}{\sqrt{n}} \|A\|_\infty \leq \|A\|_2 \leq
  \sqrt{n}\|A\|_\infty$}

  先证 $\|A\|_2 \leq \sqrt{n}\|A\|_\infty$。设 $\|x\|_2 = 1$，则
  \[
    \|Ax\|_2^2 = \sum_{i=1}^n \left| \sum_{j=1}^n a_{ij} x_j
    \right|^2 \leq \sum_{i=1}^n \left( \sum_{j=1}^n |a_{ij}| \cdot
    |x_j| \right)^2 \leq \sum_{i=1}^n \|A\|_\infty^2 \cdot \|x\|_1^2
    \leq n \|A\|_\infty^2.
  \]

  再证 $\frac{1}{\sqrt{n}} \|A\|_\infty \leq \|A\|_2$。设第 $k$ 行达到最大行和，取
  $x = (\text{sgn}(a_{k1}), \ldots, \text{sgn}(a_{kn}))^T /
  \sqrt{n}$，则 $\|x\|_2 = 1$，且
  \[
    \|Ax\|_2 \geq \frac{1}{\sqrt{n}} \left| \sum_{j=1}^n |a_{kj}|
    \right| = \frac{1}{\sqrt{n}} \|A\|_\infty.
  \]

  \subsection*{(3) $\frac{1}{\sqrt{n}} \|A\|_1 \leq \|A\|_2 \leq
  \sqrt{n}\|A\|_1$}

  由于 $\|A\|_2 = \|A^T\|_2$ 且 $\|A^T\|_\infty = \|A\|_1$，由 (2) 即得。

  \subsection*{(4) $\|A\|_2 \leq \|A\|_F \leq \sqrt{n}\|A\|_2$}

  设 $\sigma_1 \geq \sigma_2 \geq \cdots \geq \sigma_n \geq 0$ 为 $A$ 的奇异值，则
  \[
    \|A\|_2 = \sigma_1, \quad \|A\|_F = \sqrt{\sum_{i=1}^n \sigma_i^2}.
  \]

  显然 $\sigma_1^2 \leq \sum_{i=1}^n \sigma_i^2 \leq n \sigma_1^2$，即
  $\|A\|_2 \leq \|A\|_F \leq \sqrt{n}\|A\|_2$。

  \subsection*{(5) $\|AB\|_F \leq \|A\|_2\|B\|_F$}

  现设 $B = (b_1, \ldots, b_n)$ 其中 $b_j$ 为 $B$ 的第 $j$ 列。设 $C = AB$，则
  \[
    c_{ij} = \sum_{k=1}^n a_{ik} b_{kj}.
  \]
  而 $Ab_j$ 的第 $i$ 个分量为
  \[
    (Ab_j)_i = \sum_{k=1}^n a_{ik} b_{kj} = c_{ij}.
  \]
  因此
  \[
    \|Ab_j\|_2^2 = \sum_{i=1}^n (Ab_j)_i^2 = \sum_{i=1}^n c_{ij}^2,
  \]
  于是
  \[
    \|AB\|_F^2 = \sum_{i,j} c_{ij}^2 = \sum_{j=1}^n \sum_{i=1}^n
    c_{ij}^2 = \sum_{j=1}^n \|Ab_j\|_2^2.
  \]
  由相容性 $\|Ab_j\|_2 \leq \|A\|_2 \|b_j\|_2$，得
  \[
    \|AB\|_F^2 = \sum_{j=1}^n \|Ab_j\|_2^2 \leq \sum_{j=1}^n
    \|A\|_2^2 \|b_j\|_2^2 = \|A\|_2^2 \|B\|_F^2.
  \]

  \subsection*{(6) $\|AB\|_F \leq \|A\|_F\|B\|_2$}

  由 $\|AB\|_F = \|(AB)^T\|_F = \|B^T A^T\|_F$，应用 (5) 得
  \[
    \|B^T A^T\|_F \leq \|B^T\|_2 \|A^T\|_F = \|B\|_2 \|A\|_F.
  \]

\end{solution}

\section{题目 5：LU 分解的增长因子}

令 $A = LU$ 是 $n \times n$ 矩阵 $A$ 的 LU 分解且 $|l_{ij}| \leq 1$。设 $a_i^T$
和 $u_i^T$ 分别表示 $A$ 和 $U$ 的第 $i$ 行。证明
\[
  u_i^T = a_i^T - \sum_{j=1}^{i-1} l_{ij} u_j^T,
\]
以及用它证明 $\|U\|_\infty \leq 2^{n-1}\|A\|_\infty$。

\begin{solution}
  \textbf{Step 1. 证明等式 $u_i^T = a_i^T - \sum_{j=1}^{i-1} l_{ij} u_j^T$}

  由 LU 分解 $A = LU$，其中 $L$ 是下三角矩阵（对角线元素为1），$U$ 是上三角矩阵。

  考虑 $A = LU$ 的第 $i$ 行。设 $L$ 的第 $i$ 行为 $(l_{i1}, l_{i2}, \ldots,
  l_{i,i-1}, 1, 0, \ldots, 0)$，$U$ 的第 $j$ 行为 $u_j^T$。

  由矩阵乘法，$A$ 的第 $i$ 行 $a_i^T$ 为：
  \[
    a_i^T = \sum_{k=1}^{n} l_{ik} \cdot u_k^T.
  \]

  由于 $L$ 是单位下三角矩阵，$l_{ik} = 0$ 当 $k > i$，且 $l_{ii} = 1$。因此：
  \[
    a_i^T = \sum_{k=1}^{i-1} l_{ik} u_k^T + l_{ii} u_i^T =
    \sum_{k=1}^{i-1} l_{ik} u_k^T + u_i^T.
  \]

  移项即得：
  \[
    u_i^T = a_i^T - \sum_{k=1}^{i-1} l_{ik} u_k^T.
  \]

  此即所证等式。

  \textbf{Step 2. 证明 $\|U\|_\infty \leq 2^{n-1}\|A\|_\infty$}

  考虑归纳证明：
  \[
    \|u_i^T\|_1 \leq 2^{i-1} \|A\|_\infty
  \]
  为了书写方便，以下所有的范数均为 $L_1$ 范数。从而有 $\|A\|_\infty = \max_{1 \leq i \leq
  n} \|a_i^T\|$。
  于是有：

  \textbf{I.} 当 $i=1$ 时，$u_1^T = a_1^T$，因此 $\|u_1^T\| = \|a_1^T\|
  \leq \|A\|_\infty$，结论成立。

  \textbf{II.} 假设当 $i < k$ 时，$\|u_i^T\|_1 \leq 2^{i-1} \|A\|_\infty$
  成立。则当 $i = k$ 时：

  \[
    \begin{aligned}
      \|u_k^T\| & = \left\| a_k^T - \sum_{j=1}^{k-1} l_{kj} u_j^T
      \right\| \\
      & \leq \|a_k^T\| + \sum_{j=1}^{k-1} |l_{kj}| \cdot \|u_j^T\|
      \quad (\text{由三角不等式}) \\
      & \leq \|a_k^T\| + \sum_{j=1}^{k-1} 2^{j-1} \|A\|_\infty
      \quad (\text{由归纳假设和 $|l_{kj}| \leq 1$}) \\
      & \leq \|A\|_\infty + \|A\|_\infty \sum_{j=1}^{k-1} 2^{j-1} \\
      & = \|A\|_\infty \left( 1 + \sum_{j=1}^{k-1} 2^{j-1} \right) \\
      & = \|A\|_\infty \cdot 2^{k-1}.
    \end{aligned}
  \]

  从而 $\|u_k^T\|_1 \leq 2^{k-1} \|A\|_\infty$，结论成立。于是 $\|U\|_\infty =
  \max_{1 \leq i \leq n} \|u_i^T\|_1 \leq 2^{n-1} \|A\|_\infty$。
\end{solution}

\end{document}
